\section{Distance specializations and bounded transfer}
\label{sec:outer-distance-collapse}

The exact optimization also recovers the outer-distance collapse without
choosing a radius.  Indeed, every nonzero functional sector costs at least
$d(O^\perp)-1$: its label has at least $d(O^\perp)$ active blocks, the target
block may contribute zero helpers, and every other active block contributes
at least one.  The zero-functional sector attains
$M_t(D_P,K_P)+d(I^\perp)$ after minimization over $T$.  Consequently,
\begin{equation}
 d(O^\perp)\geq M_t(D_P,K_P)+d(I^\perp)+1
 \quad\Longrightarrow\quad
 \Gamma_{j,t}(O,I)=M_t(D_P,K_P)+d(I^\perp).
 \label{eq:ranked-weighted-collapse}
\end{equation}

\begin{theorem}[Confinement for a fixed target subspace]
\label{thm:objectwise-confinement}
\coverage{absent}%
Let $r$ be a nonnegative integer and let $O\leq L^N$ be an $L$-linear outer
code with $N\geq2$ and $d(O^\perp)>r+1$, and assume $0\ne T\leq W_P$.  Then
$O\circ I$ confines every
normalized recovery system for $T$ of cost at most $r$ if and only if
\begin{equation}
 r<\rho_T(I)+d(I^\perp).
 \label{eq:objectwise-confinement}
\end{equation}
Below this bound, zero-extension gives a support-preserving bijection between
the inner systems of cost at most $r$ and the corresponding systems in the
concatenation.  For an outer family with dual distance tending to infinity,
the outer-distance hypothesis eventually holds at each fixed $r$.  The
preceding equivalence then applies; the support-preserving bijection holds
exactly when~\eqref{eq:objectwise-confinement} holds.
\end{theorem}

\begin{proof}
If $0\ne B\in\mathcal B_T(O)$, choose $u\in T$ with $B(u)\ne0$.  The tuple
$B(u)\in\operatorname{FD}(O)$ has at least $d(O^\perp)$ nonzero blocks.
At most one is the target block, so every lift has at least
$d(O^\perp)-1>r$ helpers outside it.  Thus no nonzero-functional sector in
\eqref{eq:weighted-subspace-cost} occurs at cost at most $r$.

The zero-functional sector has exact cost
$\mu_{I,P,T}(0)+d(I^\perp)=\rho_T(I)+d(I^\perp)$.  Theorem
\ref{thm:ungated-ranked-confinement} now gives both directions of
\eqref{eq:objectwise-confinement} and the support-preserving bijection.
\uses{thm:ungated-ranked-confinement}%
\proves{thm:objectwise-confinement}%
\end{proof}

Minimizing the fixed-subspace threshold over all recovered subspaces of fixed
dimension gives the promised hierarchy.

\begin{theorem}[Confinement by recovered dimension under an outer-distance condition]
\label{thm:ranked-confinement}
\coverage{absent}%
Let $r$ be a nonnegative integer, let $O\leq L^N$ be $L$-linear with
$N\geq2$ and $d(O^\perp)>r+1$, and let $\ell=\dim W_P$ and
$1\leq t\leq\ell$.  The concatenation $O\circ I$ confines every
cost-$\leq r$ recovery system for every helper-recoverable
$t$-dimensional target subspace if and only if
\begin{equation}
 r<M_t(D_P,K_P)+d(I^\perp).
 \label{eq:ranked-confinement}
\end{equation}
When~\eqref{eq:ranked-confinement} holds, restriction and zero-extension
preserve every normalized equation and its exact helper support.
For an outer family with dual distance tending to infinity, the outer-distance
hypothesis eventually holds at each fixed $r$.  The preceding equivalence then
applies; preservation for all such target subspaces holds exactly
when~\eqref{eq:ranked-confinement} holds.
\end{theorem}

\begin{proof}
By Theorem~\ref{thm:relative-weight-recovery}, every such target subspace
has inner cost at least $M_t(D_P,K_P)$, and some target subspace attains this
minimum.  Apply Theorem~\ref{thm:objectwise-confinement} to each subspace.
The lower bound gives sufficiency uniformly; an attaining subspace and the
external rank-one perturbation in that proof give necessity.
\uses{thm:relative-weight-recovery, thm:objectwise-confinement}%
\proves{thm:ranked-confinement}%
\end{proof}

\subsection{One-coordinate specialization}

For $\beta\in L^*$, put
\[
 \lambda_I(\beta)=\min_{\Phi_I(y)=\beta}\wt(y),
 \qquad
 \mu_{I,x}(\beta)=
 \min_{\substack{\Phi_I(y)=\beta\\y_x\neq0}}\wt(y),
\]
with value $\infty$ when a constrained fibre is empty.  The first quantity is
the usual coset-leader weight for the inner encoder, while the second retains
the target-coordinate constraint.  The functional-dual decomposition itself
is classical concatenation machinery
\cite[Theorem~2.1]{ChenLingXing2005}.
Here a dual word ``through $x$'' simply means that its coefficient at $x$ is
nonzero.  The quantity below counts total dual-word weight; the general
recovery cost excludes the forced target coordinate, so for a nonzero target
column the two conventions satisfy $Z_{j,x}=\Gamma_{j,T}+1$.

\begin{theorem}[Exact weighted finite-length confinement at one coordinate]
\label{thm:weighted-pointed-confinement}
\coverage{absent}%
Fix an outer block $j$, an inner coordinate $x$, and $N\geq2$, and assume
that the coordinate projection $O\to L$ at block $j$ is nonzero, hence
surjective.  The minimum total Hamming weight of a word in $(O\circ I)^\perp$
that is nonzero at $(j,x)$ and is not the zero-extension of a word of
$I^\perp$ in block $j$ is
\begin{equation}
 Z_{j,x}(O,I)=
 \min\left\{
  \mu_{I,x}(0)+d(I^\perp),
  \min_{\substack{\beta=(\beta_h)\in\operatorname{FD}(O)\\\beta\neq0}}
  \left(\mu_{I,x}(\beta_j)+
        \sum_{h\neq j}\lambda_I(\beta_h)\right)
 \right\}.
 \label{eq:weighted-pointed-cost}
\end{equation}
A minimum over an empty set is $\infty$.  Consequently, for every $r\geq0$,
restriction to block $j$ and zero-extension are inverse, support-preserving
bijections between the normalized recovery equations through $x$ in $I$ and
those through $(j,x)$ in $O\circ I$, both using at most $r$ helpers, if and
only if
\begin{equation}
 Z_{j,x}(O,I)>r+1.
 \label{eq:weighted-pointed-gate}
\end{equation}
The same equivalence holds for the exact helper-support families.  When these
conditions hold, the inclusion-minimal recovery sets are copied as well.
\end{theorem}

\begin{proof}
Write a concatenated dual word in blocks and apply
\eqref{eq:block-functional-dual}.  For a fixed nonzero functional tuple, the
block representatives minimize independently, giving the second term of
\eqref{eq:weighted-pointed-cost}.  For the zero tuple, nonconfinement costs
$\mu_{I,x}(0)$ in the target block and $d(I^\perp)$ in another block.  The
projection hypothesis makes confined words exactly the zero-extensions of
inner-dual words, so these two sectors prove~\eqref{eq:weighted-pointed-cost}.

If the target column is nonzero, take $P=\{x\}$ and let $T\leq U_P$ be its
span.  After choosing a nonzero vector of $T$, a linear map $T\to L^*$ is one
functional.  Normalizing its target coefficient and rescaling the functional
tuple identify the minima in~\eqref{eq:weighted-subspace-cost} with those in
\eqref{eq:weighted-pointed-cost}.  The target coordinate contributes one to
total word weight and is not a helper, so
\[
 Z_{j,x}(O,I)=\Gamma_{j,T}(O,I)+1.
\]
Theorem~\ref{thm:ungated-ranked-confinement} gives the claimed condition.  If the
target column is zero, the same condition follows directly from the exact minimum
already proved.  In either case, below the minimum every equation is confined,
so restriction and zero-extension are inverse.  In one dimension every exact
helper support is the support of its normalized equation; the equivalence also
holds for exact helper-support families.  It implies equality of their
inclusion-minimal members, but minimal-support equality need not imply
equation confinement, as Theorem~\ref{thm:minimal-support-confinement} shows.
\uses{thm:ungated-ranked-confinement}%
\proves{thm:weighted-pointed-confinement}%
\end{proof}

Theorem~\ref{thm:weighted-pointed-confinement} retains information that the
outer support distance discards.  The following family makes the difference
strict without using a geometric inner code.  Its construction chooses an
outer dual repetition line whose projective multipliers avoid every ratio
between low-cost inner labels.  Consequently the labelled criterion holds even
though ordinary distance sees no margin.

\begin{proposition}[A family from a Singer cycle without the outer-distance condition]
\label{prop:strict-weighted-transfer}
\coverage{absent}%
\imports{Singer:projective-cycle}%
Let $k\geq2$ and suppose
\begin{equation}
 \frac{q^k-1}{q-1}>(k+1)^2.
 \label{eq:singer-disjointness}
\end{equation}
Let $I$ be the $[k+1,k,2]_q$ single-parity-check code generated by
\[
 G=(I_k\mid\mathbf1),
\]
and identify its message space with $L=\F_{q^k}$.  For every
$k+1\leq n\leq2k+1$, there is an $L$-linear length-$n$ outer code $O$ with
$d(O^\perp)=n$ such that~\eqref{eq:weighted-pointed-gate} holds at helper
radius $n-1$ for every inner coordinate and every outer block.  Thus all
normalized radius-$(n-1)$ recovery equations and their exact helper supports
are copied blockwise although the usual outer-distance condition fails: here
$d(O^\perp)=n$ is not strictly greater than $r+1=n$.

In particular, this holds for the $[4,3,2]_5$ inner code and a length-five
outer code at helper radius four.
\end{proposition}

\begin{proof}
The $k+1$ coordinate functionals of the encoder give a set $S$ of $k+1$
points in the projective space on $L^*$.  The multiplication maps of $L$,
modulo $\F_q^*$, induce on that projective space a Singer cycle of order
$(q^k-1)/(q-1)$, acting regularly on its points~\cite{Singer1938}.
Indeed, under the trace identification, the multiplier class sending a
nonzero functional class $[f]$ to $[g]$ is uniquely $[g/f]$.
Each ordered pair of points of $S$ therefore forbids at most one cycle
element, so at most $|S|^2=(k+1)^2$ elements send a point of $S$ to a point
of $S$.  Thus~\eqref{eq:singer-disjointness} permits a multiplication map
$a:L\to L$, explicitly $a(u)=\gamma u$ for some $\gamma\in L^\times$, whose
dual action satisfies
$a^*(S)\cap S=\varnothing$.

Set
\[
 O=\{u\in L^n:u_0+a(u_1)+u_2+\cdots+u_{n-1}=0\}.
\]
Every functional-dual tuple has the form
$(f,f\circ a,f,\ldots,f)$.  If $f\neq0$, all $n$ entries are nonzero, so
$d(O^\perp)=n$.  A realization of total weight $n$ would have weight one in
every block.  Weight-one realizations are exactly the scalar multiples of
the $k+1$ coordinate functionals, so both $[f]$ and $[f\circ a]$ would lie
in $S$, contrary to the choice of $a$.  Every nonzero tuple therefore costs
at least $n+1$.

Every outer coordinate projection is surjective.  Finally, $I^\perp$ is
spanned by $(1,\ldots,1,-1)$, so
$d(I^\perp)=\mu_{I,x}(0)=k+1$ for every $x$.  The zero-functional cost in
\eqref{eq:weighted-pointed-cost} is $2k+2>n$, and the nonzero-functional cost
is at least $n+1$.  Thus $Z_{j,x}(O,I)>n$ for every $j,x$.  At $r=n-1$ the
ordinary distance condition would require $d(O^\perp)>n$, which is false, while the exact
weighted condition holds.  Taking $(q,k,n)=(5,3,5)$ gives the stated concrete
instance.  Since the inner recovery equations have $k$ helpers and
$n-1\geq k$, the asserted transfer is nonvacuous throughout the family.
\uses{thm:weighted-pointed-confinement}%
\proves{prop:strict-weighted-transfer}%
\end{proof}

\paragraph{Explicit field and multiplier.}
For $(q,k,n)=(5,3,5)$, take
$L=\F_5[\alpha]/(\alpha^3+\alpha+1)$ in the basis $(1,\alpha,\alpha^2)$.
The cubic has no root in $\F_5$, so it is irreducible.
Multiplication by $\gamma=\alpha+\alpha^2$ has matrix
\[
 M_\gamma=\begin{pmatrix}0&4&4\\1&4&3\\1&1&4\end{pmatrix}.
\]
The coordinate functionals of $G=(I_3\mid\mathbf1)$ are the rows
$e_1^*,e_2^*,e_3^*,(1,1,1)$.  Their pullbacks by $M_\gamma$ are
\[
 (0,4,4),\quad(1,4,3),\quad(1,1,4),\quad(2,4,1).
\]
None is projectively a coordinate row or $(1,1,1)$, proving the required
disjointness directly.  The outer code is
$O=\{z\in L^5:z_0+\gamma z_1+z_2+z_3+z_4=0\}$.
This specifies the encoder alignment as well as the field and parameters.

For a nonzero, helper-recoverable singleton target $P=\{x\}$, the first relative weight is
the least number of helpers in a dual word through $x$.  If the total-weight
threshold used for one-coordinate equations is denoted by $z_x(I)$, then
\begin{equation}
 z_x(I)=M_1(D_P,K_P)+1+d(I^\perp).
 \label{eq:singleton-threshold}
\end{equation}
The helper-radius condition~\eqref{eq:ranked-confinement} is therefore
$r+1<z_x(I)$, with no change of convention hidden in the reduction.
If the target column is nonzero but not helper-recoverable, then $W_P=0$:
its local recovery cost is infinite and there is no first relative weight
of this pair.
Theorem~\ref{thm:ungated-ranked-confinement} is the common finite formula.
The outer-distance condition removes its nonzero-functional sector and gives the
RGHW threshold in every recovered dimension.  Restricting instead to one
target coordinate gives Theorem~\ref{thm:weighted-pointed-confinement} without
that condition.

The labelled functions suffice to compose fibre minima.  Computing
nonconfinement also retains the nonzero-kernel cost $d(I^\perp)$ through
\eqref{eq:dual-distance-composition}; its ordinary zero-label minimum is
zero and does not encode this cost.  Minimizing lifts must also be stored
to propagate witnesses.  Coarser summaries answer different questions.
Section~\ref{sec:separations} shows that
relative-weight minima do not determine recovery reliability and that the
abstract nested pair does not determine confinement.
